%!TEX program = xelatex
\documentclass[11pt]{article}

\usepackage[margin=1in]{geometry}
\usepackage{amsmath, amssymb, amsthm}
\usepackage{graphicx}
\usepackage{hyperref}
\usepackage{xcolor}
\usepackage{booktabs}
\usepackage{multirow}
\usepackage{float}
\usepackage[section]{placeins}
\usepackage{caption}
\usepackage{adjustbox}

\newcommand{\safeincludegraphics}[2][]{%
  \IfFileExists{#2}{\includegraphics[#1]{#2}}{%
    \fbox{\parbox[c][0.22\textheight][c]{0.9\linewidth}{\centering [Placeholder for Figure: \texttt{\detokenize{#2}}]}}%
  }%
}
\newcommand{\figw}{0.62\linewidth}
\newcommand{\tablesize}{\footnotesize}
\newcommand{\idea}[1]{\textcolor{red}{[#1]}}
\newcommand{\question}[1]{\textcolor{blue}{[#1]}}


\setlength{\textfloatsep}{10pt plus 2pt minus 2pt}
\setlength{\floatsep}{8pt plus 2pt minus 2pt}
\captionsetup{font=small,skip=4pt}

\newtheorem{theorem}{Theorem}
\newtheorem{assumption}{Assumption}
\newtheorem{proposition}{Proposition}
\newtheorem{corollary}{Corollary}

\setlength{\parindent}{0pt}
\setlength{\parskip}{1\baselineskip}

\title{Evolution of MoE Routing: From Low-Rank Projections to Debiased Variance-Aware (Second-Order) Routing}
\author{Theory Group}
\date{\today}

\begin{document}
\maketitle

% {\color{red}
% \section{Comments}

Your current discussion mixes several different operations that actually serve different purposes. The key point is that mean removal, spike suppression, and tail lifting (e.g., whitening) are three different operations with three different objectives, and they should not be conflated.

Subtracting the mean only removes the rank-1 DC offset of the representation. This is useful in cases where the mean acts as a shared non-informative carrier—for example in MoE routing or low-bit quantization, where a large common shift wastes dynamic range or dominates linear projections. In such cases, centering improves the signal-to-noise ratio without changing the overall geometry. However, mean removal does not solve the broader anisotropy problem, because the spectrum can still remain highly skewed: a few dominant singular directions may still concentrate most of the energy while the long tail remains weak. Therefore, mean removal should be viewed mainly as a bias correction, not a general solution to representation anisotropy.

Whitening is a much stronger operation. In addition to removing the mean, it rescales the covariance so that different directions have comparable variance, effectively lifting the long-tail directions relative to dominant principal modes. This makes whitening a more natural tool if the goal is to treat anisotropy in a general sense. However, whitening can also distort useful structure, amplify noise, and destroy meaningful non-isotropic organization in the representation space. In many cases the dominant directions may contain useful shared features rather than purely nuisance bias.

Because of this, the “mean issue” is not a single problem with a single solution. The effect of the mean depends on the downstream mechanism. In quantization the mean wastes representational range, so separating it helps. In routing the mean may dominate a linear selector, so centering or spike attenuation may help. But in optimization or representation learning the mean may be partly harmful and partly useful because it contributes both to curvature concentration and to shared feature organization.

Therefore, the right framework should distinguish analysis space from intervention space. Analytically, you can decompose the representation into the mean component, dominant singular directions, and the long tail. But intervention should be task-specific: centering for quantization or routing, low-rank suppression for curvature concentration, and only cautious whitening when full spectral rebalancing is truly necessary.

In other words, mean removal is a local surgical tool, while whitening is a global geometric transform. If the goal is to build a general theory, the object of study should not be the mean alone but the broader spike–tail geometry and how different system components interact with it. 
% }
% \newpage

\section{Recent Updates}

% \idea{The opening should now be organized around three clear points: (1) the theoretical progress on centered quadratic routing, (2) what the current mean-related experiments do and do not imply, and (3) why ACP matters for the story. Each update should be explicit, compact, and connected to the continual-learning objective.}

Our current main question is whether \emph{centered, debiased quadratic routing} is actually necessary for continual-learning-oriented routing, or whether a centered linear router already captures the relevant structure. In this version, \texttt{linear\_running\_center} is treated as the strict baseline, while ACP-style models are treated as structured mean-direction comparators rather than fair baselines.

\begin{enumerate}
    \item Established the new core design principle of the quadratic router.

    \textbf{Conclusion:} the router input must be centered at both first and second order. First-order centering removes the dominant mean-driven initialization imbalance, and second-order centering removes the raw-quadratic positive bias that otherwise creates initialization-time load skew. This gives a principled candidate mechanism for quadratic routing, but not yet a license to treat it as the established mainline; current train-time evidence for quadratic routing remains limited.

    \item Clarified what the current mean-related experiments do and do not show.

    \textbf{Conclusion:} the current running-center, low-rank, and ACP variants show that our expert-quality metrics do not map one-to-one to downstream performance. In particular, stronger compactness or isolation does not automatically imply the best task accuracy. Therefore the continual-learning story must be written in terms of routing-induced geometry and forgetting-related structure, not in terms of a single scalar metric.

    \item Identified what the ACP result actually means for the story.

    \textbf{Conclusion:} ACP shows that the mean direction is not simply useless. The right conclusion is subtler: raw mean dominance is harmful because it causes collapse, but once the center direction is isolated and handled structurally, it may still carry useful routing information. However, ACP-style models should be interpreted as parameter-augmented comparators rather than strict baselines, because they add extra pre-logit capacity. This makes a lighter comparator such as ACP-lite methodologically important.

\end{enumerate}

\subsection{Problem setting and thoughts.}

\idea{Now we should set our goals. The overall goal we want to achieve now is 'Continual Learning'. To achieve continual learning, we think we must through gating make our expert specialized, thus we have to let our router have the ability that in the spectral space to separate the data, now we first think of using our variation to distribute. 
But actually, we don't know if the linear matrix is unable to use spectral feature to distribute. And we need to define 'Distribute according to the variation properly.'
Thus in this iteration, in addition to evaluating whether the routing has this abiliity, we should also develop and define the problem deeper.
}

\subsection{ToDo:}
\begin{enumerate}
    \item Give a direct theoretical and empirical answer to the key question for quadratic routing: after subtracting the second-order mean, when does the router still preserve enough variance signal to separate domains or semantic groups?
    \item Test whether quadratic routing actually produces more separated and lower-rank expert covariances than centered linear routing, using both no-training audits and trained-model geometry analysis.
    \item Tighten the link from routing to continual learning by showing that better second-order partitioning reduces expert overlap and can plausibly improve the forgetting proxy.
    \item Clarify the relation between quadratic routing and expert low-rankness: determine whether variance-aware routing is itself sufficient to drive lower-rank expert representations, or whether additional structural constraints are required.
\end{enumerate}

\section{Background and Theoretical Critique}

% \idea{In this part, we need to clarify what we are sure of and what our chain of thoughts are. The problem}


The current evidence forces a narrower story than before. We no longer just need a router that is expressive. We need one explanation that simultaneously accounts for four facts: centering fixes balance, low-rank routing can improve geometry, ACP-style mean-direction handling can improve some downstream results, and uncentered routing can look confident while remaining geometrically poor.

Two issues therefore become central. The first is \emph{functional form}: after centering, the remaining domain signal may live mainly in covariance rather than in mean, so a purely linear router becomes structurally misspecified. The second is \emph{geometric measurement}: even if a router compresses expert inputs, the usual $ED_{\mathrm{comp}}$ can still overcount weak spectral tails, so ``low-rank'' itself needs a stricter definition.

This is why router analysis and expert analysis cannot be separated. The router determines which activations are pooled together inside each expert, hence it determines the covariance, effective dimension, and spectral interference that the expert must absorb.

\subsection{Continual Learning, expert low-rankness, and routing}
Let $C_{e,D}$ denote the covariance matrix of the activations routed to expert $e$ from domain $D$. The central CL intuition is:
\begin{enumerate}
    \item the router decides which samples enter expert $e$,
    \item the induced partition determines $C_{e,D}$,
    \item a low-rank or low-effective-dimension covariance means the occupied representation subspace is small,
    \item if the subspaces of old and new domains overlap less, catastrophic forgetting is weaker.
\end{enumerate}
This suggests the following qualitative forgetting proxy for expert $e$:
\[
\mathcal{F}_e(D_{\mathrm{old}}\!\to\! D_{\mathrm{new}})
\propto
p(e\mid D_{\mathrm{new}})
\cdot
\left\|U_{e,\mathrm{old}}^\top U_{e,\mathrm{new}}\right\|_F^2,
\]
where $U_{e,\mathrm{old}}$ and $U_{e,\mathrm{new}}$ are principal subspaces extracted from the two domain-specific covariances. The first factor is controlled by router exclusivity; the second is controlled by geometric compactness and cross-domain isolation.

This leads to a sharper recognition than before: the router is not just choosing experts, it is \emph{sculpting covariance}. Centering is therefore a causal intervention on which statistic the router can see. Before centering, the router mostly reacts to the global mean direction; after centering, it can start reacting to covariance structure.

\subsection{The Curse of Mean Bias and why centering matters}
Recent observations suggest that language model activations contain a strong shared mean bias. Write
\[
x=\mu+\epsilon,\qquad \mathbb{E}[\epsilon]=0,\qquad \|\mu\|_2\gg \|\epsilon\|_2 \ \text{at initialization scale}.
\]
Consider a randomly initialized router with logits $\ell_e = w_e^\top x$, where each router row satisfies $w_e\sim \mathcal{N}(0,\sigma_w^2 I)$.

\begin{theorem}[Centering removes the dominant Step-0 variance source]
Conditioning on the input vector $x$, the random logit $\ell_e=w_e^\top x$ satisfies
\[
\mathrm{Var}_{w_e}(\ell_e\mid x)=\sigma_w^2\|x\|_2^2.
\]
If we center by the true mean and set $x_c=x-\mu=\epsilon$, then
\[
\mathrm{Var}_{w_e}(w_e^\top x_c\mid x)=\sigma_w^2\|\epsilon\|_2^2.
\]
\end{theorem}

\begin{proof}
Condition on a fixed realization of $x$. Because $w_e\sim\mathcal{N}(0,\sigma_w^2 I)$, we have
\[
\mathbb{E}_{w_e}[w_e^\top x\mid x]
=
\mathbb{E}_{w_e}[w_e]^\top x
=
0.
\]
Hence
\[
\mathrm{Var}_{w_e}(w_e^\top x\mid x)
=
\mathbb{E}_{w_e}\big[(w_e^\top x)^2\mid x\big].
\]
Using $\mathbb{E}[w_e w_e^\top]=\sigma_w^2 I$,
\[
\mathbb{E}_{w_e}\big[(w_e^\top x)^2\mid x\big]
=
x^\top \mathbb{E}_{w_e}[w_e w_e^\top] x
=
x^\top (\sigma_w^2 I)x
=
\sigma_w^2\|x\|_2^2.
\]
This proves the first identity. Since $x=\mu+\epsilon$ and $x_c=x-\mu=\epsilon$, the same calculation with $x_c$ in place of $x$ gives
\[
\mathrm{Var}_{w_e}(w_e^\top x_c\mid x)
=
\sigma_w^2\|x_c\|_2^2
=
\sigma_w^2\|\epsilon\|_2^2.
\]
\end{proof}

If $\|\mu\|_2\gg\|\epsilon\|_2$, then the initial logit variance is dominated by the shared mean direction. This gives the first precise effect of centering: it does \emph{not} merely improve training convenience; it changes the pre-training logit distribution itself.

\subsection{The limits of linear routing after centering}
Centering solves the mean-bias problem, but not the full routing problem. If two domains share near-zero centered means but differ strongly in covariance, then a linear router can only exploit residual first-order structure. The key point is not that linear routing has literally zero signal in every such setting; the key point is that its function class is generally mismatched to the Bayes-optimal rule.

\subsection{A proof path: covariance-only domain shift}
To move beyond intuition, consider the minimal stylized setting
\[
x\mid D=A \sim \mathcal{N}(0,\Sigma_A),
\qquad
x\mid D=B \sim \mathcal{N}(0,\Sigma_B),
\]
with equal class priors and diagonal covariances in dimension $d$. Throughout this subsection, the domain label is $y\in\{A,B\}$, all expectations are population expectations, and the diagonal covariance restriction is an explicit modeling assumption rather than a universal claim. This is the cleanest model where first-order information vanishes while second-order structure remains.

\begin{theorem}[Centered first-order scores have zero domain mean gap]
For any vector $a$,
\[
\mathbb{E}[a^\top x \mid D=A] - \mathbb{E}[a^\top x \mid D=B] = 0.
\]
\end{theorem}

\begin{proof}
By linearity of expectation,
\[
\mathbb{E}[a^\top x\mid D=A]
=
a^\top \mathbb{E}[x\mid D=A],
\qquad
\mathbb{E}[a^\top x\mid D=B]
=
a^\top \mathbb{E}[x\mid D=B].
\]
Under the model assumptions, both conditional means are zero. Therefore
\[
\mathbb{E}[a^\top x \mid D=A] - \mathbb{E}[a^\top x \mid D=B]
=
a^\top 0-a^\top 0
=
0.
\]
\end{proof}

\begin{proposition}[Zero-threshold linear rules are uninformative under centered symmetric domains]
Suppose $x\mid D=A$ and $x\mid D=B$ are both centrally symmetric about the origin, meaning that $x$ and $-x$ have the same conditional law in each domain. Then for every nonzero vector $a$,
\[
\mathbb P(a^\top x>0\mid D=A)
=
\mathbb P(a^\top x>0\mid D=B)
=
\frac12.
\]
\end{proposition}

\begin{proof}
Fix a domain $D\in\{A,B\}$ and define the scalar random variable $s=a^\top x$. Since $x$ and $-x$ have the same conditional law, $s$ and $-s$ also have the same conditional law. Hence
\[
\mathbb P(s>0\mid D)=\mathbb P(-s>0\mid D)=\mathbb P(s<0\mid D).
\]
For a non-degenerate Gaussian or any absolutely continuous centrally symmetric law, $\mathbb P(s=0\mid D)=0$, so
\[
\mathbb P(s>0\mid D)+\mathbb P(s<0\mid D)=1.
\]
Therefore $\mathbb P(s>0\mid D)=1/2$. Since the argument is identical for $D=A$ and $D=B$, the claim follows.
\end{proof}

This proposition isolates a narrower claim than ``linear routing is blind'': without an explicit threshold shift, centered symmetric domains give no zero-threshold linear separation at all.

If one allows an affine score $s(x)=a^\top x+b$, then the situation changes. In the Gaussian model,
\[
\mathbb P(s(x)>0\mid y)
=
\Phi\!\left(\frac{b}{\sqrt{a^\top\Sigma_y a}}\right),
\]
so thresholded linear rules can retain partial signal through variance-dependent tail probabilities. This is why the correct statement is \emph{misspecification}, not literal blindness: the linear function class may extract some signal, but it still cannot represent the Bayes-optimal quadratic boundary in the covariance-only model.

Now define
\[
\Delta = \mathrm{diag}(\Sigma_A)-\mathrm{diag}(\Sigma_B),
\qquad
\lambda = \tfrac{1}{2}\big(\mathrm{diag}(\Sigma_A)+\mathrm{diag}(\Sigma_B)\big),
\]
and $q_{\mathrm{deb}} = x\odot x - \lambda$.

\begin{theorem}[Debiased quadratic features recover covariance separation]
If $b=\Delta$, then
\[
\mathbb{E}[b^\top q_{\mathrm{deb}}\mid D=A]
-
\mathbb{E}[b^\top q_{\mathrm{deb}}\mid D=B]
=
\|\Delta\|_2^2.
\]
\end{theorem}

\begin{proof}
For each coordinate $i$ and each domain $D\in\{A,B\}$, the Gaussian is zero mean, so
\[
\mathbb{E}[x_i^2\mid D]
=
\mathrm{Var}(x_i\mid D)
=
\Sigma_{D,ii}.
\]
Therefore
\[
\mathbb{E}[x\odot x\mid D=A]=\mathrm{diag}(\Sigma_A),
\qquad
\mathbb{E}[x\odot x\mid D=B]=\mathrm{diag}(\Sigma_B).
\]
Subtracting $\lambda=\tfrac12(\mathrm{diag}(\Sigma_A)+\mathrm{diag}(\Sigma_B))$ yields
\[
\mathbb{E}[q_{\mathrm{deb}}\mid D=A]=\tfrac{1}{2}\Delta,
\qquad
\mathbb{E}[q_{\mathrm{deb}}\mid D=B]=-\tfrac{1}{2}\Delta.
\]
Hence
\[
\mathbb{E}[b^\top q_{\mathrm{deb}}\mid D=A]
-
\mathbb{E}[b^\top q_{\mathrm{deb}}\mid D=B]
=
b^\top\Big(\tfrac12\Delta-(-\tfrac12\Delta)\Big)
=
b^\top \Delta.
\]
Setting $b=\Delta$ gives
\[
b^\top \Delta
=
\Delta^\top \Delta
=
\|\Delta\|_2^2.
\]
\end{proof}

Theorem 3 only establishes nonzero second-order separation in the debiased feature space; it does not claim that $b = \Delta$ coincides with the Bayes-optimal coefficient vector.

\begin{theorem}[Bayes-optimal routing is diagonal quadratic in the Gaussian covariance-only model]
Under the same assumptions,
\[
\log \frac{p(x\mid D=A)}{p(x\mid D=B)}
=
c-\tfrac{1}{2}x^\top(\Sigma_A^{-1}-\Sigma_B^{-1})x,
\]
for some constant $c$. When $\Sigma_A,\Sigma_B$ are diagonal, this rule is affine in $x\odot x$, equivalently affine in $q_{\mathrm{deb}}$.
\end{theorem}

\begin{proof}
The zero-mean Gaussian density with covariance $\Sigma_D$ is
\[
p(x\mid D)
=
\frac{1}{(2\pi)^{d/2}\det(\Sigma_D)^{1/2}}
\exp\!\left(-\tfrac12 x^\top \Sigma_D^{-1}x\right).
\]
Taking the log-likelihood ratio gives
\[
\log \frac{p(x\mid D=A)}{p(x\mid D=B)}
=
-\tfrac12 \log\det(\Sigma_A)
+\tfrac12 \log\det(\Sigma_B)
-\tfrac12 x^\top \Sigma_A^{-1}x
+\tfrac12 x^\top \Sigma_B^{-1}x.
\]
Collecting the constant terms into
\[
c=-\tfrac12 \log\det(\Sigma_A)+\tfrac12 \log\det(\Sigma_B),
\]
we obtain
\[
\log \frac{p(x\mid D=A)}{p(x\mid D=B)}
=
c-\tfrac12 x^\top(\Sigma_A^{-1}-\Sigma_B^{-1})x.
\]
If $\Sigma_A$ and $\Sigma_B$ are diagonal, then $\Sigma_A^{-1}-\Sigma_B^{-1}$ is also diagonal. Writing
\[
\Sigma_A^{-1}-\Sigma_B^{-1}=\mathrm{diag}(\gamma_1,\dots,\gamma_d),
\]
the quadratic form becomes
\[
x^\top(\Sigma_A^{-1}-\Sigma_B^{-1})x
=
\sum_{i=1}^d \gamma_i x_i^2,
\]
which is affine in $x\odot x$. Since $x\odot x=q_{\mathrm{deb}}+\lambda$, the same rule is also affine in $q_{\mathrm{deb}}$, with the shift $\lambda$ absorbed into the constant term.
\end{proof}

This gives the correct logical separation. Theorem 2 only removes first-moment separation. Theorem 4 is what upgrades the claim: in the covariance-only Gaussian regime, centered linear routing is generally misspecified, while debiased diagonal quadratic routing has the correct functional form.

\subsection{Why only diagonal second-order information?}
The proposed router is not full covariance-aware; it is deliberately \emph{diagonal second-order}. This choice has three roles.
\begin{enumerate}
    \item \textbf{Minimal functional upgrade.} A full quadratic router of the form $x^\top M_e x$ requires $\mathcal O(Ed^2)$ parameters and computation, whereas a diagonal term $b_e^\top (x\odot x)$ keeps both at $\mathcal O(Ed)$.
    \item \textbf{Feasible online debiasing.} A full running estimate of $\mathbb E[xx^\top]$ requires $\mathcal O(d^2)$ state and is substantially noisier. A diagonal running estimate of $\mathbb E[x\odot x]$ is cheap and naturally compatible with EMA updates.
    \item \textbf{Methodological scope.} The goal of this work is to test whether the cheapest second-order correction already changes routing geometry in a CL-relevant way.
\end{enumerate}

\begin{proposition}[Diagonal quadratic features do not see purely off-diagonal covariance shifts]
Suppose two zero-mean domains satisfy
\[
\mathrm{diag}(\Sigma_A)=\mathrm{diag}(\Sigma_B).
\]
Then for every vector $b$,
\[
\mathbb E[b^\top (x\odot x)\mid D=A]
=
\mathbb E[b^\top (x\odot x)\mid D=B].
\]
\end{proposition}

\begin{proof}
For any domain $D$ and coordinate $i$,
\[
\mathbb E[x_i^2\mid D]=\Sigma_{D,ii}.
\]
Hence
\[
\mathbb E[b^\top(x\odot x)\mid D]
=
\sum_{i=1}^d b_i \Sigma_{D,ii}
=
b^\top \mathrm{diag}(\Sigma_D).
\]
If the two domains share the same diagonal covariance, the two expectations are equal.
\end{proof}

This proposition states the method boundary explicitly: the router only exploits diagonal second-order structure. If two domains differ only through off-diagonal covariance, the proposed feature map cannot separate them at the level of first moments.

\subsection{Why raw quadratic features also collapse}
Now let $x_c$ be centered and define the raw diagonal quadratic feature
\[
q_{\mathrm{raw}} = x_c \odot x_c.
\]

\begin{theorem}[Raw quadratic features regenerate a positive mean]
If $\mathbb{E}[x_c]=0$ and $\mathrm{Cov}(x_c)=\Sigma$, then
\[
\mathbb{E}[q_{\mathrm{raw}}]
=
\mathbb{E}[x_c\odot x_c]
=
\mathrm{diag}(\Sigma)
\triangleq \lambda
\succeq
\mathbf{0}.
\]
\end{theorem}

\begin{proof}
For each coordinate $i$,
\[
\mathbb{E}[(x_c)_i^2]
=
\mathrm{Var}((x_c)_i)+\big(\mathbb{E}[(x_c)_i]\big)^2.
\]
Because $\mathbb{E}[x_c]=0$, the second term vanishes, so
\[
\mathbb{E}[(x_c)_i^2]
=
\mathrm{Var}((x_c)_i)
=
\Sigma_{ii}.
\]
Stacking these coordinate-wise identities yields
\[
\mathbb{E}[x_c\odot x_c]
=
\mathrm{diag}(\Sigma)
\triangleq \lambda.
\]
Since every variance is nonnegative, $\lambda\succeq 0$ componentwise. If every coordinate variance is strictly positive, then $\lambda$ is strictly positive componentwise.
\end{proof}

Thus raw quadratic routing reintroduces a deterministic positive shift. The stronger initialization statement is the following.

\begin{proposition}[Initialization-time load skew for a raw quadratic router]
Let $z\sim \mathcal N(0,\Lambda)$ with $\Lambda=\mathrm{diag}(\nu_1,\dots,\nu_r)$, and consider a raw quadratic router
\[
\ell_e(z)=a_e^\top z + q_e^\top (z\odot z),
\qquad
q_{e,i}\overset{\mathrm{i.i.d.}}{\sim}\mathcal N\!\left(0,\frac{\sigma_q^2}{r}\right).
\]
Define $\lambda=(\nu_1,\dots,\nu_r)^\top$ and
\[
m_e=\mathbb E_z[\ell_e(z)\mid \theta].
\]
Then
\[
m_e=q_e^\top \lambda,
\qquad
\mathrm{Var}_\theta(m_e)=\frac{\sigma_q^2}{r}\|\lambda\|_2^2.
\]
Moreover, if the logits remain in a small-initialization regime where the softmax admits a first-order expansion with quadratic remainder, then the conditional expected load satisfies
\[
L_e
=
\mathbb E_z[\mathrm{softmax}_e(\ell(z))\mid\theta]
=
\frac1E+\frac1E(m_e-\bar m)+\mathcal O\!\big(\mathbb E_z[\|\ell(z)\|_2^2\mid\theta]\big),
\]
where $\bar m=\frac1E\sum_{j=1}^E m_j$.
\end{proposition}

\begin{proof}
Since $z$ is centered, $\mathbb E[z]=0$, so
\[
\mathbb E_z[a_e^\top z\mid\theta]=0.
\]
Also, for each coordinate $i$,
\[
\mathbb E_z[z_i^2\mid\theta]=\nu_i.
\]
Therefore
\[
m_e
=
\mathbb E_z[\ell_e(z)\mid\theta]
=
\sum_{i=1}^r q_{e,i}\nu_i
=
q_e^\top \lambda.
\]
Because the $q_{e,i}$ are independent with variance $\sigma_q^2/r$,
\[
\mathrm{Var}_\theta(m_e)
=
\sum_{i=1}^r \nu_i^2 \mathrm{Var}(q_{e,i})
=
\frac{\sigma_q^2}{r}\sum_{i=1}^r \nu_i^2
=
\frac{\sigma_q^2}{r}\|\lambda\|_2^2.
\]

For the expected load, let $s_e(\ell)=\exp(\ell_e)/\sum_{j=1}^E \exp(\ell_j)$. A first-order Taylor expansion of $s_e$ around $\ell=0$ gives
\[
s_e(\ell)
=
\frac1E + \frac1E\Big(\ell_e-\frac1E\sum_{j=1}^E \ell_j\Big) + R_e(\ell),
\]
where the remainder satisfies $|R_e(\ell)|\le C\|\ell\|_2^2$ in a neighborhood of the origin. Taking conditional expectation with respect to $z$ yields
\[
L_e
=
\frac1E+\frac1E\Big(m_e-\bar m\Big)+\mathbb E_z[R_e(\ell(z))\mid\theta].
\]
The bound on the remainder gives
\[
\mathbb E_z[R_e(\ell(z))\mid\theta]
=
\mathcal O\!\big(\mathbb E_z[\|\ell(z)\|_2^2\mid\theta]\big),
\]
which proves the claim.
\end{proof}

Raw quadratic routing therefore regenerates a deterministic second-order mean term. Under anisotropic Gaussian residuals and a small-logit softmax expansion, this induces initialization-time expert-dependent logit mean shifts and hence non-uniform expected loads. This is the precise initialization-level sense in which raw quadratic routing exhibits a second collapse mechanism. The result is stated for expected softmax loads under small-logit expansion; discrete top-k routing requires a separate approximation argument.

\subsection{Router geometry, centering, and the empirical contradictions we must explain}
Our 24L experiments already show that:
\begin{itemize}
    \item \texttt{low\_rank} achieves the lowest $ED_{\mathrm{comp}}$ and lowest spectral interference (SI), indicating strong geometric compactness and isolation;
    \item \texttt{linear\_running\_center} is the most balanced downstream model, leading in 5/7 tasks;
    \item \texttt{linear\_no\_center} attains perfect perturbation consistency and the largest margin, yet also yields the worst SI, revealing that confidence or stability alone does not imply useful specialization;
    \item \texttt{acp\_on} becomes strongest in some later-stage SI and downstream comparisons, suggesting that centering direction itself carries structural information that can be exploited more carefully.
\end{itemize}
These observations force a richer story: centering, low-rankness, and router specificity are causally linked, but not reducible to a single scalar notion such as ``confidence'', ``stability'', or ``effective dimension''.

\section{Hypotheses and Mathematical Modeling}

\idea{According to our fixed insights, we need our representation x to subtract a running mean, and also the squared term to subtract a running mean(of the squared term), the exact form needs modifying.}

Based on the critique above, we update the working assumptions and make the modeling explicit.

\begin{assumption}[Variance-driven specialization after centering]
After removing the shared macroscopic mean drift, domain-specific semantics are primarily expressed by covariance or energy structure rather than by first-order mean offsets. A router that wishes to induce expert specialization should therefore combine first-order centering with second-order debiasing.
\end{assumption}

\subsection{Methodology: the debiased second-order router}
We propose the \emph{Debiased Centered Diagonal Quadratic Router}. The running first-order mean is tracked by
\[
\mu_t = (1-\beta_t)\mu_{t-1} + \beta_t \bar{x}_t,
\qquad
x_c = x-\mu_t.
\]
The raw quadratic feature is
\[
q_{\mathrm{raw}} = x_c \odot x_c.
\]
Its running mean is tracked by
\[
\lambda_t = (1-\beta_t)\lambda_{t-1} + \beta_t \overline{q_{\mathrm{raw}}}_t,
\qquad
q_{\mathrm{deb}} = q_{\mathrm{raw}}-\lambda_t.
\]
The final expert logit becomes
\[
\ell_e(x)=a_e^\top x_c + b_e^\top q_{\mathrm{deb}}.
\]
Under stationarity this makes the residual means small in expectation, but the statement is approximate under drift because the router uses online EMA estimates rather than oracle moments.

\subsection{Methodological note: ACP as a structured mean-direction comparator}
ACP is not a second-order router. It is a comparison method showing that the center direction may still contain useful information once it is isolated rather than left entangled with the raw feature stream. It should not be treated as a strict baseline in fairness-sensitive comparisons, because adding a dedicated pre-logit decomposition also adds extra expressive capacity. Let
\[
u_t=\frac{\mu_t}{\|\mu_t\|_2},
\qquad
x_c=x-\mu_t.
\]
ACP decomposes the centered feature into the component parallel to the center direction and the orthogonal residual:
\[
x_{\parallel}=(u_t^\top x_c)u_t,
\qquad
x_{\perp}=x_c-x_{\parallel}.
\]
The routing feature is then constructed as
\[
\tilde x = W_{\parallel}x_{\parallel}+W_{\perp}x_{\perp},
\]
followed by an ordinary linear gate
\[
\ell_e^{\mathrm{ACP}}(x)=g_e^\top \tilde x.
\]
Its conceptual role in this report is important: ACP shows that the mean direction should not be treated as universally useless. The correct interpretation is that \emph{raw mean dominance} is harmful because it causes imbalance, whereas an explicitly separated center-direction component may still provide useful routing structure. For the next iteration, a lighter comparator (\emph{ACP-lite}) is methodologically preferable: split the feature into mean-direction projection and centered residual, then apply two matched $E\times d$ linear maps whose logits are summed. This tests the mean-direction hypothesis with a fairer parameter budget.

\begin{proposition}[EMA tracking under stationarity]
Let
\[
\hat v_t=(1-\beta)\hat v_{t-1}+\beta q_t,
\qquad
\mathbb E[q_t]=v^\star,
\]
with constant $\beta\in(0,1]$ and stationary input sequence $(q_t)$. Then
\[
\mathbb E[\hat v_t-v^\star]=(1-\beta)^t(\hat v_0-v^\star).
\]
\end{proposition}

\begin{proof}
Subtract $v^\star$ from both sides:
\[
\hat v_t-v^\star
=
(1-\beta)(\hat v_{t-1}-v^\star)+\beta(q_t-v^\star).
\]
Taking expectation and using $\mathbb E[q_t-v^\star]=0$ gives
\[
\mathbb E[\hat v_t-v^\star]
=
(1-\beta)\mathbb E[\hat v_{t-1}-v^\star].
\]
Iterating this recursion yields
\[
\mathbb E[\hat v_t-v^\star]=(1-\beta)^t(\hat v_0-v^\star).
\]
\end{proof}

Therefore EMA debiasing is exact only in an oracle sense. In practice it is an approximation whose error depends on both stochastic noise and distribution drift. In particular, for any router coefficient vector $b$,
\[
\left|b^\top(q_{\mathrm{raw}}-\hat v_t)-b^\top(q_{\mathrm{raw}}-v^\star)\right|
\le
\|b\|_1\|\hat v_t-v^\star\|_\infty,
\]
so moment-tracking error transfers linearly into logit error. Under non-stationary training dynamics, the relevant decomposition is a lag--variance tradeoff rather than an exact equality.

\subsection{Why expert specificity must be multi-dimensional}
The current empirical evidence strongly suggests that expert specificity cannot be defined by one number. The existing Round1--3 scripts already decompose it into distinct dimensions:
\begin{itemize}
    \item \textbf{Assignment exclusivity}: $H(E\mid C)$ and, in future, NMI$(E,C)$ or purity;
    \item \textbf{Geometric compactness}: $ED$ and $ED_{\mathrm{comp}}$;
    \item \textbf{Isolation}: Overlap@r and SI;
    \item \textbf{Stability / confidence}: CR$(\sigma)$ and router margin;
    \item \textbf{Fairness}: $CV_{\mathrm{load}}$.
\end{itemize}
Each captures a different failure mode. For example, a collapsed router may have low entropy and high perturbation consistency, yet still be useless for CL because it mixes domains inside one expert or leaves no room for balanced specialization.

\subsection{Metric-design hypothesis and script-grounded definitions}
\textbf{Hypothesis.} Expert specificity should be defined as a multi-view structural state rather than as a single score.

\textbf{Mathematical modeling.} The existing Round1--3 scripts already instantiate a concrete measurement interface:
\[
H(E\mid C)
=
-\sum_c p(c)\sum_e p(e\mid c)\log_2 p(e\mid c),
\]
\[
ED(C)=\frac{(\sum_i \nu_i)^2}{\sum_i \nu_i^2},
\qquad
Overlap@r(i,j)=\frac{\|U_i^\top U_j\|_F^2}{r},
\]
\[
CV_{\mathrm{load}}
=
\frac{\sqrt{\frac{1}{E}\sum_{e=1}^E (p_e-\frac{1}{E})^2}}{\frac{1}{E}},
\qquad
CR(\sigma)=\frac{1}{N}\sum_{n=1}^N \mathbf{1}[e_{\mathrm{clean}}^{(n)}=e_{\mathrm{noisy},\sigma}^{(n)}],
\]
\[
SI_e=\frac{\mathrm{Tr}(C_{A,e}C_{B,e})}{\|C_{A,e}\|_F\|C_{B,e}\|_F},
\qquad
Margin=\ell_{(1)}-\ell_{(2)}.
\]
Here $ED_{\mathrm{comp}}$ denotes the imbalance-compensated effective dimension already used in Round1. Together these quantities form the structural specificity object
\[
\mathcal{S}
=
\big(H(E\mid C),\ ED_{\mathrm{comp}},\ SI,\ CV_{\mathrm{load}};\ CR,\ Margin\big).
\]

\textbf{Experiment design.} We should continue reporting these six quantities layerwise and at model mean, then add normalized assignment metrics such as NMI$(E,C)$ and Purity$(C\mid E)$ so that comparisons across different expert counts and cluster granularities remain meaningful.

\textbf{Conclusion / Update.} The report should standardize a four-part panel, but the low-rank slot itself still needs refinement; $ED_{\mathrm{comp}}$ is a good current baseline, not necessarily the final scientific answer.

\subsection{Which low-rank metric should be primary?}
Let $\nu_1,\dots,\nu_d$ be the eigenvalues of the expert covariance $C$, and define $p_i=\nu_i/\mathrm{Tr}(C)$. The current metric
\[
ED(C)=\frac{(\sum_i \nu_i)^2}{\sum_i \nu_i^2}
=
\frac{1}{\sum_i p_i^2}
=
\exp(H_2(p))
\]
is the participation ratio, i.e.\ the exponential of the order-2 R\'enyi entropy. It is scale-invariant and useful for cross-model comparison, but it still counts weak spectral tails.

There are at least four serious candidates:
\[
\mathrm{erank}(C)=\exp\!\Big(-\sum_i p_i\log p_i\Big),
\qquad
\mathrm{sr}(C)=\frac{\mathrm{Tr}(C)}{\lambda_{\max}(C)},
\]
\[
d_\tau(C)=\mathrm{Tr}\!\big(C(C+\tau I)^{-1}\big)=\sum_i \frac{\nu_i}{\nu_i+\tau},
\qquad
r_\alpha(C)=\min\Big\{k:\frac{\sum_{i\le k}\nu_i}{\sum_j \nu_j}\ge \alpha\Big\}.
\]

\begin{theorem}[Statistical dimension is a soft count of usable directions]
Let $C\succeq 0$ have eigenvalues $\nu_1,\dots,\nu_d$, and let $\tau>0$. Define $N_{>\tau}(C)=|\{i:\nu_i>\tau\}|$. Then
\[
\frac{1}{2}N_{>\tau}(C)
\le
d_\tau(C)
\le
N_{>\tau}(C) + \frac{1}{\tau}\sum_{\nu_i\le \tau}\nu_i.
\]
\end{theorem}

\begin{proof}
Write
\[
d_\tau(C)=\sum_{i=1}^d \frac{\nu_i}{\nu_i+\tau},
\qquad \tau>0.
\]
Partition the index set into
\[
I_> = \{i:\nu_i>\tau\},
\qquad
I_{\le} = \{i:\nu_i\le\tau\}.
\]
If $i\in I_>$, then $\nu_i>\tau>0$, hence
\[
\frac{\nu_i}{\nu_i+\tau}>\frac{1}{2}
\qquad\text{and}\qquad
\frac{\nu_i}{\nu_i+\tau}<1.
\]
If $i\in I_{\le}$, then $\nu_i\ge 0$ and
\[
0\le \frac{\nu_i}{\nu_i+\tau}\le \frac{\nu_i}{\tau},
\]
because $\nu_i+\tau\ge\tau$.

For the lower bound, keep only the contribution of $I_>$:
\[
d_\tau(C)
=
\sum_{i\in I_>} \frac{\nu_i}{\nu_i+\tau}
\;+\;
\sum_{i\in I_{\le}} \frac{\nu_i}{\nu_i+\tau}
\ge
\sum_{i\in I_>} \frac{1}{2}
=
\frac{1}{2}|I_>|
=
\frac{1}{2}N_{>\tau}(C).
\]
For the upper bound, use the estimates above:
\[
d_\tau(C)
\le
\sum_{i\in I_>} 1
\;+\;
\sum_{i\in I_{\le}} \frac{\nu_i}{\tau}
=
|I_>|
\;+\;
\frac{1}{\tau}\sum_{i\in I_{\le}}\nu_i
=
N_{>\tau}(C)+\frac{1}{\tau}\sum_{\nu_i\le\tau}\nu_i.
\]
This proves the claim.
\end{proof}

This theorem makes $d_\tau$ especially attractive in our setting: it behaves like a soft count of directions that rise above the noise floor $\tau$, whereas $ED$ and $\mathrm{erank}$ are more affected by long weak tails, and stable rank is overly dominated by the top eigen-direction.

\subsection{A principled metric panel}
We therefore treat the following four quantities as the \emph{primary specificity panel}:
\[
H(E\mid C),\qquad
\mathcal{R}_{\mathrm{geom}},\qquad
SI,\qquad
CV_{\mathrm{load}}.
\]
Here $\mathcal{R}_{\mathrm{geom}}$ denotes the geometry slot. For the current report it is filled by $ED_{\mathrm{comp}}$ for backward comparability; for the next iteration, the main upgrade candidate is the noise-aware statistic $d_{\tau,\mathrm{comp}}$, with $r_{0.9,\mathrm{comp}}$ as a human-interpretable side metric.

They answer four different questions:
\begin{enumerate}
    \item does the router assign similar semantic clusters to a concentrated set of experts?
    \item are the expert inputs geometrically compact?
    \item do distinct domains interfere inside the same expert?
    \item is the resulting specialization real, or merely a consequence of collapse?
\end{enumerate}
We regard CR and margin as \emph{auxiliary} metrics because they measure robustness and decisiveness, but not necessarily CL-friendly specialization. Accordingly, a model should not be summarized by a single scalar headline; it should be judged by non-domination on the primary panel.

For the next iteration, $\tau$ must also be chosen by a fixed protocol rather than by visual convenience. The default proposal is to define $\tau_0$ as the median of the bottom 20\% eigenvalues pooled across experts at a given layer, report $d_{\tau_0,\mathrm{comp}}$ as the main geometry metric, and add a sensitivity check at $\{0.5\tau_0,\tau_0,2\tau_0\}$.

\subsection{An exploratory composite score}
If a single summary number is eventually needed, an exploratory index can be defined as
\[
S_{\mathrm{assign}} = 1-\frac{H(E\mid C)}{\log E},
\quad
S_{\mathrm{geom}} = 1-\frac{ED_{\mathrm{comp}}}{D},
\]
\[
S_{\mathrm{iso}} = 1-SI,
\quad
S_{\mathrm{bal}} = e^{-CV_{\mathrm{load}}},
\quad
S_{\mathrm{stab}} = \mathrm{AUC}_{\sigma}(\mathrm{CR}(\sigma)),
\]
\[
ESS =
\left(
S_{\mathrm{assign}}
S_{\mathrm{geom}}
S_{\mathrm{iso}}
S_{\mathrm{bal}}
S_{\mathrm{stab}}
\right)^{1/5}.
\]
However, we do \emph{not} recommend using $ESS$ as the primary headline metric yet. The normalization choices are still somewhat subjective, and the current counterexamples (especially \texttt{linear\_no\_center}) show that a single-score collapse can hide mechanistically important distinctions.

\section{Empirical Study}

% \idea{Here we need to give a recap of all the experiments we have done and each experiment's firm conclusion and what we haven't tested, what we are about to test. Each subsection must be neatly organized into a complete iteration content.}

\subsection{Empirical recap across completed iterations}
Table~\ref{tab:24l-summary} summarizes the most informative 24L checkpoint-40000 comparison.

\begin{table}[H]
\centering
\tablesize
\begin{adjustbox}{max width=\linewidth}
\begin{tabular}{lcccc}
\toprule
Model & $ED_{\mathrm{comp}} \downarrow$ & CR@$\sigma{=}0.5 \uparrow$ & Mean SI $\downarrow$ & Downstream Avg $\uparrow$ \\
\midrule
\texttt{linear\_running\_center} & 66.85 & 0.9500 & 0.6763 & 0.4537 \\
\texttt{low\_rank} & \textbf{64.12} & 0.9430 & \textbf{0.6389} & 0.4510 \\
\texttt{acp\_on} & 66.47 & 0.9496 & 0.6895 & \textbf{0.4586} \\
\texttt{linear\_no\_center} & 63.79 & \textbf{1.0000} & 0.7952 & 0.4478 \\
\bottomrule
\end{tabular}
\end{adjustbox}
\caption{Cross-round summary on the 24L ckpt40000 comparison. The table already reveals that structural compactness, robustness, interference, and downstream quality do not collapse to a single ordering.}
\label{tab:24l-summary}
\end{table}

\paragraph{Conclusion.} The current evidence is internally consistent with a multi-view notion of specialization: \texttt{linear\_running\_center} is the strict baseline, \texttt{low\_rank} is structurally informative, ACP-style models may improve some downstream or SI-related outcomes but should be read as comparators rather than fair baselines, and \texttt{linear\_no\_center} is a crucial counterexample showing that large margin and perfect CR can still coexist with poor spectral isolation.

\subsection{Round1: structural specialization}
Round1 evaluates expert geometry using $ED$, $ED_{\mathrm{comp}}$, Overlap@r, $H(E\mid C)$, and $CV_{\mathrm{load}}$. The key lesson is that low-rank routing indeed compresses expert input geometry: in the 24L setting, \texttt{low\_rank} achieves the best $ED_{\mathrm{comp}}=64.12$, while \texttt{linear\_running\_center} and \texttt{acp\_on} remain in the mid-60s.

\paragraph{Conclusion.} Low-rank routing strongly supports the geometric compactness part of the story, but Round1 alone cannot determine whether the compactness is CL-friendly. That requires cross-domain isolation and robustness measurements.

\subsection{Round2: perturbation stability}
Round2 measures routing consistency under Gaussian perturbations. It is tempting to read high CR as a proxy for good specialization, but the results prove that this interpretation is incomplete: \texttt{linear\_no\_center} attains CR$=1.0$ for all noise levels, yet it does not produce the best geometry or the best cross-domain isolation.

\paragraph{Conclusion.} Stability is a necessary descriptive property of the router, but not a sufficient criterion for useful expert specificity. In particular, perfect CR can reflect a degenerate routing regime rather than a desirable one.

\subsection{Round3: spectral interference and margin}
Round3 directly addresses the CL-relevant geometric question by measuring SI and margin. In the 24L comparison, \texttt{low\_rank} obtains the lowest SI (0.6389), which is consistent with the idea that compressed expert geometry reduces cross-domain overlap. At the same time, \texttt{linear\_no\_center} obtains the largest margin but also the worst SI, proving that confidence is not identical to isolation.

The ACP observation remains useful, but it must now be interpreted more carefully: when an ACP-style model improves SI or downstream quality, this tells us that mean-direction structure may carry useful routing information after disentanglement. It does \emph{not} by itself prove that ACP is a better baseline, because the comparison is not parameter-neutral.

\paragraph{Conclusion.} SI is currently the closest metric to the CL mechanism we care about. Margin remains important, but only when interpreted jointly with SI and load balance.

\paragraph{Scope note on quadratic evidence.} At the moment, trained-model evidence for quadratic routing is materially weaker than for the centered linear family because quadratic training has only produced usable measurements at ckpt40000. Therefore any report claim about quadratic routing should be written as a candidate route supported by partial evidence, not as a checkpoint-stable conclusion across training.

\subsection{Interpretation: what centering actually changes}
The centering story must therefore be stated in three layers:
\begin{enumerate}
    \item \textbf{Optimization layer}: centering removes the dominant initialization-time variance source and is therefore expected to reduce Step-0 load skew;
    \item \textbf{Representation layer}: centering removes the shared first-order drift and exposes the residual structure where domain variation lives;
    \item \textbf{Geometric layer}: once the nuisance mean is removed, the router can be evaluated on whether it induces compact and isolated expert subspaces rather than merely reflecting a global bias direction.
\end{enumerate}
This explains why \texttt{linear\_running\_center} is such a strong baseline: even before introducing second-order statistics, centering already removes a major confounder and yields a balanced, task-competitive router.

\subsection{Next sub-iteration: no-training experiments}
\idea{This corresponds to the new minimal discriminative test designed in v5.1 iteration state.}
The next empirical package should deliberately avoid further training and operate on frozen dense or MoE activations. Each check is meant to isolate one theoretical claim before optimization dynamics are allowed to confound the result.
\begin{enumerate}
    \item \textbf{Dense mean-bias audit.} \textbf{Purpose:} test the claim that the dominant Step-0 variance source is the shared mean direction. \textbf{Steps:} measure the layerwise ratio $\|\mu^\star\|_2^2/\mathbb E\|x\|_2^2$, then simulate random linear routers on raw and centered activations and compare load CV and logit variance. \textbf{Expected conclusion:} centering sharply reduces initialization variance and improves Step-0 balance whenever the mean-bias ratio is large.
    \item \textbf{First-order vs second-order separability audit.} \textbf{Purpose:} test whether covariance signal remains after first-moment centering. \textbf{Steps:} on held-out domain pairs, compare plug-in linear scores $a^\top x_c$ and plug-in debiased quadratic scores $b^\top q_{\mathrm{deb}}$ without training, using ROC/AUC or balanced accuracy as the readout. \textbf{Expected conclusion:} when domains are mainly separated by diagonal energy structure, the debiased quadratic score outperforms the centered linear score.
    \item \textbf{Step-0 raw-vs-debiased quadratic audit.} \textbf{Purpose:} test the initialization-time load-skew theorem for raw quadratic routing. \textbf{Steps:} compare \texttt{raw\_linear}, \texttt{centered\_linear}, \texttt{raw\_quadratic}, and \texttt{debiased\_quadratic} on the same frozen activations with random weights; record mean logits, load CV, and the correlation between empirical load skew and the predicted shift $m_e=q_e^\top \lambda$. \textbf{Expected conclusion:} \texttt{raw\_quadratic} shows expert-dependent mean shifts and larger load imbalance, while \texttt{debiased\_quadratic} suppresses this effect.
    \item \textbf{EMA residual-mean audit.} \textbf{Purpose:} test the approximation gap between oracle debiasing and online EMA debiasing. \textbf{Steps:} measure $\|\mathbb E[\hat x_{c,t}]\|$ and $\|\mathbb E[q_{\mathrm{deb}}]\|$ globally and per domain under several $\beta$ values, then compare these residuals with the induced logit error bound $\|b\|_1\|\hat v_t-v^\star\|_\infty$. \textbf{Expected conclusion:} residual means decay under stationarity but expose a clear lag--noise tradeoff under drift.
    \item \textbf{Frozen partition audit.} \textbf{Purpose:} test whether analytic first-order and second-order partitions already induce different expert geometry. \textbf{Steps:} use argmax partitions induced by the plug-in scores above and evaluate $ED_{\mathrm{comp}}$, Overlap, SI, and margin on those partitions directly. \textbf{Expected conclusion:} if the second-order advantage is intrinsic rather than optimization-induced, the geometry gap should appear before training.
    \item \textbf{Noise-aware rank audit.} \textbf{Purpose:} test whether the current low-rank conclusions are robust to weak spectral tails. \textbf{Steps:} compute $d_{\tau_0,\mathrm{comp}}$ and $r_{0.9,\mathrm{comp}}$ alongside $ED_{\mathrm{comp}}$, with a sensitivity table at $\{0.5\tau_0,\tau_0,2\tau_0\}$. \textbf{Expected conclusion:} if a model is genuinely compact rather than merely tail-heavy, its geometry ranking should remain stable across this protocol.
\end{enumerate}

\paragraph{Conclusion.} These no-training experiments are not auxiliary conveniences; they are the cleanest way to separate representational geometry from optimization dynamics and to decide whether the advantage of second-order routing is intrinsic or merely training-induced.

\section{Next iteration}

The next iteration has four tightly coupled goals:
\begin{enumerate}
    \item \textbf{Verify the centering mechanism.} Run the dense mean-bias audit and the EMA residual-mean audit on frozen dense and MoE representations. The goal is to determine whether centering and second-order debiasing are actually removing the intended nuisance moments, and to identify the useful $\beta$ regime.
    \item \textbf{Verify the second-order initialization claim.} Run the Step-0 raw-vs-debiased quadratic audit and compare the empirical load skew against the predicted initialization shift. The goal is to test the theorem at the load level rather than only at the feature-mean level.
    \item \textbf{Introduce a fair mean-direction comparator.} Evaluate ACP-lite at ckpt40000 so that ``explicitly use the mean direction'' can be tested without the confound of a larger pre-logit parameter budget.
    \item \textbf{Standardize the specificity panel.} Report \{$H(E\mid C)$, $\mathcal{R}_{\mathrm{geom}}$, SI, $CV_{\mathrm{load}}$\} as the primary panel, keep CR and margin as auxiliary descriptors, and evaluate whether $d_{\tau,\mathrm{comp}}$ with the fixed $\tau$ protocol should replace $ED_{\mathrm{comp}}$ as the main geometry metric.
    \item \textbf{Tighten the CL connection.} Use the frozen partition audit and the updated geometry panel to connect routing decisions to expert covariance overlap and then to a more explicit CL-style forgetting proxy in the subsequent section.
\end{enumerate}

The main scientific objective is now sharper than before: we are no longer merely asking whether a router is sparse or low-rank, but whether centered linear routing already suffices, whether mean-direction structure helps under fair comparison, and whether debiased second-order routing creates the \emph{right kind} of expert geometry for continual learning.

\end{document}
